\section{Introduction}
\label{sec:introduction}

Agentic AI systems are evolving from passive assistants into autonomous decision-makers that plan, call tools, and act on retrieved data~\cite{mialon2023gaia,drouin2024workarena,xie2024osworld,xi2024agentir}. Their failures extend beyond incorrect answers to goal drift, unsafe tool use, and socially harmful behaviour: a recent coding agent deleted a live production database under an explicit code freeze~\cite{nolan2025replit}, and indirect prompt injection through retrieved documents silently redirects tool-augmented agents~\cite{greshake2023indirect,zhou2024poisonedrag}. Trustworthiness must therefore be understood as a deployment-level concern shaped by realistic environmental, operational, and social conditions.

Despite this shift, current evaluation practices remain limited in two compounding ways. First, ``trustworthiness'' is used inconsistently---as safety, helpfulness, or an umbrella term. Existing surveys list many desiderata~\cite{liu2024trustllm,wang2024trustworthyagent} but offer no small, measurable property set against which an agent can be assessed. Second, existing benchmarks focus on isolated slices---task-solving~\cite{zhou2023webarena}, coding~\cite{jimenez2023swebench}, hallucination~\cite{li2023halueval}, jailbreak resistance~\cite{chao2024jailbreakbench}, or tool use~\cite{li2023apibank}---constructed around narrow tasks and single failure modes~\cite{liang2023helm,kiela2021dynabench}, so high scores do not imply real-world reliability. The central limitation is thus the lack of \emph{both} a concrete trustworthiness definition and a principled notion of \emph{representativeness}: agent trustworthiness is a distributional property, and misalignment between benchmark and deployment distributions creates false confidence~\cite{koh2021wilds,shirali2022dynamicbench}, especially for low-frequency, high-consequence events central to trustworthy deployment.

To address both gaps, we propose the \textbf{Holographic Agent Assessment Framework (HAAF)}. The key idea is to first \emph{define} agent trustworthiness as a five-property profile (P1--P5, \S\ref{subsec:defining_trustworthiness}) and then \emph{measure} it over a \emph{scenario manifold} spanning task types, tool interfaces, interaction dynamics, social contexts, and risk levels, rather than over disconnected benchmark islands. A \emph{distribution-aware representative sampling engine} defines \emph{what} to test, while three probing interfaces---static cognitive and policy analysis, interactive sandbox simulation, and social-ethical alignment---define \emph{how}; the present instantiation (\S\ref{sec:experiments}) exercises the interactive sandbox layer in depth, with the other layers reserved for future cycles. The components are connected through an iterative \emph{Trustworthy Optimization Factory}: red-team probing exposes vulnerabilities, blue-team hardening designs targeted interventions, and re-evaluation verifies improvement. Rather than evaluating model weights in isolation, we assess \emph{deployed agentic systems} comprising model, prompting policies, tool interfaces, guardrails, and oversight.

\paragraph{Representativeness: target vs.\ current evidence.}
We separate the \emph{target} of representativeness from the \emph{evidence} we offer. Our 100 hand-authored scenarios split into 24 design-study cases (s01--s24) driving the focal-model Factory cycle and 76 validation cases (s25--s100) instantiating the five-property profile. The five sandbox tools (\texttt{search\_docs}, \texttt{db\_query}, \texttt{read\_file}, \texttt{write\_file}, \texttt{send\_message}) each map to a top-$K$ operation class in public-trace agent benchmarks~\cite{jimenez2023swebench,li2023apibank,zhou2023webarena,liu2023agentbench}. We claim the \emph{target} and \emph{methodology} (Layer-4 sampling objective with a deployment-grounded weighting hook), but \emph{not} that the present scenarios constitute a statistically representative slice of production traffic; grounding weights in real telemetry is committed as v2 (\S\ref{sec:conclusion}).

\paragraph{Contributions.}
This paper makes six contributions:
\begin{enumerate}[leftmargin=*,label=(C\arabic*)]
\item \textbf{Definition.} An operational five-property definition (P1--P5) grounded in current AI risk frameworks~\cite{nist2023airmf,eu2024aiact}, with \emph{Improper Refusal} (IR) as a first-class taxonomy class mapped to P1 rather than a side metric.
\item \textbf{Cross-model evaluation.} 13 contemporary models across seven families (Llama, Mistral, Kimi, GLM, Qwen, GPT-oss, DeepSeek) on a 100-scenario suite, producing $2{,}600$ Control/Treated trajectories---broader than the 24-scenario design study.
\item \textbf{End-to-end instantiation.} A single shared sandbox with red-team attribution, blue-team interventions, and re-evaluation: 12 of 13 systems improve on $\mathrm{RWF}$ without per-model tuning; Qwen3-32B is intervention-resistant ($\Delta{=}+0.019$).
\item \textbf{Layer-4 reference implementation.} A reference implementation of the distribution-aware sampling objective; a 40-of-100 subset preserves the cross-model ranking (Kendall-$\tau$\,=\,0.890, Spearman-$\rho$\,=\,0.963).
\item \textbf{Statistical confidence.} Wilson 95\%~CIs on per-property success rates and paired-scenario bootstrap ($B{=}10{,}000$) on four intra-family anti-scaling pairs; two reject the larger-is-no-worse null at $p<0.05$, a third at the boundary ($p=0.050$). With $n{=}14$--$26$ scenarios per property, per-cell CIs span $\pm 0.2$, so these are descriptive of cross-family trends rather than powered per-property claims (\S\ref{sec:discussion}).
\item \textbf{Explicit scope.} Scenarios and tools are hand-authored, IR validation uses a single LLM-judge (Kimi-K2.5, itself one of the 13 evaluated agents---a deliberate scope limitation discussed in \S\ref{sec:discussion}), and scenario weights are not yet grounded in production traces. A deployment-grounded sampler is committed as v2.
\end{enumerate}
